\SourceLeaf{349}%
die gekrümmte Fläche durch eine Gleichung zwischen \(x,y,z\) bestimmt ist. Noch komplizierter wird er, wenn die Fläche dadurch gegeben ist, dass \(x,y,z\) als Funktionen zweier neuer Variablen \(p,q\) dargestellt sind. Im letzten Fall enthält der Ausdruck fünfzehn Elemente, nämlich die partiellen Ableitungen erster und zweiter Ordnung von \(x,y,z\) nach \(p\) und \(q\). Seine Bedeutung liegt vor allem darin, dass er zu einem anderen Ausdruck führt, der zu den bemerkenswertesten Sätzen dieser Theorie gehört. Bei dieser Darstellung der gekrümmten Fläche hat der allgemeine Ausdruck für ein beliebiges Linienelement auf ihr,
\[
\text{oder für}\quad\BareRoot(\mathrm{d}x^2+\mathrm{d}y^2+\mathrm{d}z^2),\quad\text{die Form}\quad\BareRoot(E\mathrm{d}x^2+2F\mathrm{d}x\mathord{.}\mathrm{d}y+G\mathrm{d}y^2)
\]
wobei \(E,F.\ G\) wiederum Funktionen von \(p\) und \(q\) sind. Der erwähnte neue Ausdruck für das Krümmungsmaß enthält nur diese Größen und ihre partiellen Ableitungen erster und zweiter Ordnung. Zur Bestimmung des Krümmungsmaßes genügt also der allgemeine Ausdruck für ein Linienelement, ohne dass die Ausdrücke für die Koordinaten \(x,y,z\) selbst bekannt sein müssen. Daraus folgt unmittelbar der bemerkenswerte Lehrsatz, dass beim Abwickeln einer gekrümmten Fläche oder eines Stücks davon auf eine andere Fläche das Krümmungsmaß in jedem Punkt unverändert bleibt. Als Sonderfall folgt daraus, dass auf einer gekrümmten Fläche, die sich in eine Ebene abwickeln lässt, das Krümmungsmaß überall \(=0\) ist. Daraus ergibt sich unmittelbar die charakteristische Gleichung der in eine Ebene abwickelbaren Flächen, wenn man \(z\) als Funktion von \(x\) und \(y\) auffasst,
\[
\frac{\mathrm{dd}z}{\mathrm{d}x^2}\mathord{.}\frac{\mathrm{dd}z}{\mathrm{d}y^2}-(\frac{\mathrm{dd}z}{\mathrm{d}x\mathord{.}\mathrm{d}y})^2=0
\]
eine Gleichung, die zwar schon lange bekannt ist, nach dem Urteil des Verfassers aber bisher ohne den erforderlichen strengen Beweis geblieben war.

Diese Sätze eröffnen einen neuen Blick auf die Theorie der gekrümmten Flächen und damit ein weites, noch ganz unbebautes Untersuchungsfeld. Betrachtet man Flächen als Körper mit einer verschwindenden Dimension, lässt ihre Rolle als Grenzen von Körpern beiseite und fasst sie zugleich als biegsam und undehnbar auf, so erkennt man, dass zwei wesentlich verschiedene Arten von
